% Distortion (jfet-distortion): the set-up a person would wire, semantically the deck's.
%   VDD vdd 0 18, RD vdd d <rd>       the supply and the drain load                  -> source V_DD (set) + resistor R_D
%   RS s 0 <rs>                       the source resistor, unbypassed                -> resistor R_S
%   RIN in g 10000, VIN in 0 SIN(…)   the 1 kHz sine through 10 kΩ                   -> resistor R_in + source v_in (set)
%   J1 d g s DUT, V(d) read           copy A, the common-source stage; its output    -> njfet, voltmeter v_out (measured)
%   J1 vdd g s DUT, V(s) read         copy B, the follower deck: drain on the rail   -> voltmeter v_out on the source
%   The follower is a second deck; the gain-finding probe run is not drawn.
\documentclass[tikz, border=10pt]{standalone}
\usepackage{bench}
\begin{document}
\begin{tikzpicture}
\benchtitle{0,7.8}{Distortion}
\def\gnd{2.0}
% the transistor
\node[njfet] (QA) at (4,3) {};
\node[anchor=west, font=\small, align=left] at ($(QA.D)+(0.15,-0.55)$) {A};
% copy A: the drain load up to the rail
\draw (QA.D) -- ++(0,0.4) coordinate (Adj) to[R, l_={$R_D$}] ++(0,1.6) coordinate (Atop);
\coordinate (Avtop) at ($(Atop)+(4.2,0)$);
\draw (Atop) -- (Avtop);
\draw[set] (Avtop) to[V, l={$V_{DD}$}] (Avtop |- 0,\gnd) node[ground, black]{};
\nodename{anchor=south}{$(Atop)+(1.0,0.05)$}{vdd}
\nodename{anchor=east}{$(QA.D)+(-0.15,0.2)$}{d}
% the gate: the source resistance and the signal source
\coordinate (Ag) at (QA.G);
\draw (Ag) -- ++(-0.3,0) coordinate (Agj) to[R, l_={$R_{in}$}] ++(-1.8,0) coordinate (Avin);
\draw[set] (Avin) to[V, l_=$v_{in}$] (Avin |- 0,\gnd) node[ground, black]{};
\nodename{anchor=south}{$(Ag)+(-0.15,0.05)$}{g}
\nodename{anchor=south east}{$(Avin)+(-0.05,0.05)$}{in}
% the source resistor to ground
\draw (QA.S) to[R, l={$R_S$}] (QA.S |- 0,\gnd) node[ground]{};
\nodename{anchor=east}{$(QA.S)+(-0.1,-0.2)$}{s}
% the voltmeter on Adj
\coordinate (Avo) at ($(Adj)+(2.6,0)$);
\draw (Adj) -- (Avo);
\voltmeter{Avo |- 0,\gnd}{Avo}{$v_{out}$}
\draw (Avo |- 0,\gnd) node[ground]{};
% the transistor
\node[njfet] (QB) at (14,3) {};
\node[anchor=west, font=\small, align=left] at ($(QB.D)+(0.15,-0.55)$) {B};
% copy B: the drain on the rail
\draw (QB.D) -- ++(0,2.0) coordinate (Btop);
\coordinate (Bvtop) at ($(Btop)+(4.2,0)$);
\draw (Btop) -- (Bvtop);
\draw[set] (Bvtop) to[V, l={$V_{DD}$}] (Bvtop |- 0,\gnd) node[ground, black]{};
\nodename{anchor=south}{$(Btop)+(1.0,0.05)$}{vdd}
% the gate: the source resistance and the signal source
\coordinate (Bg) at (QB.G);
\draw (Bg) -- ++(-0.3,0) coordinate (Bgj) to[R, l_={$R_{in}$}] ++(-1.8,0) coordinate (Bvin);
\draw[set] (Bvin) to[V, l_=$v_{in}$] (Bvin |- 0,\gnd) node[ground, black]{};
\nodename{anchor=south}{$(Bg)+(-0.15,0.05)$}{g}
\nodename{anchor=south east}{$(Bvin)+(-0.05,0.05)$}{in}
% the source resistor to ground
\draw (QB.S) to[R, l_={$R_S$}] (QB.S |- 0,\gnd) node[ground]{};
\nodename{anchor=west}{$(QB.S)+(0.1,-0.3)$}{s}
\coordinate (Bsj) at ($(QB.S)+(0,-0.5)$);
% the voltmeter on Bsj
\coordinate (Bvo) at ($(Bsj)+(2.2,0)$);
\draw (Bsj) -- (Bvo);
\voltmeter{Bvo |- 0,\gnd}{Bvo}{$v_{out}$}
\draw (Bvo |- 0,\gnd) node[ground]{};
% the definitions
\definitions{0,-0.1}{%
  $V_{DD} = 18$ V;\quad $R_S$ for $I_D = I_{DSS}/3$ (from the transfer sweep), unbypassed;\quad $R_D$ for $V_{DS} = 9$ V;\quad $R_{in} = 10$ k\ohm\\[2pt]
  $v_{in}$ = 1 kHz sine, its level set by a 20 mV probe run for 1 V peak at the output (and 3 V for A, into the gate diode)\\[2pt]
  A: common source, the output on the drain;\quad B: source follower, the drain on the rail, the output across $R_S$\\[2pt]
  read $v_{out}$ over the last four of ten cycles: THD, HD2, HD3 by Goertzel%ENDDEFS
}
\end{tikzpicture}
\end{document}
